\subsection{Trazabilidad de consulta de actividad}\label{anexo:bitacora}
Las \cref{sec:diseno-bitacora,sec:impl-bitacora,sec:pruebas-bitacora} separan
la consulta filtrada de la verificación criptográfica y del alcance completo
de las \cref{tab:cu-17,tab:rf-20}. La evidencia se atribuye a cada grupo ejecutado;
la presencia de una prueba o una ruta de código no acredita su ejecución global.

Los veinte casos de aplicación se encuentran en
\rutarepo{crates/application/tests/audit_query/contract.rs} y
\rutarepo{crates/application/tests/audit_query/service.rs}. Ejercitan entradas,
cursor, autorización, límites y respuestas del puerto. El uso de dobles de
almacenamiento en ese nivel no representa una ejecución de PostgreSQL.

{\emergencystretch=3em Las catorce comprobaciones reales del adaptador se distribuyen en tres módulos
del directorio \rutarepo{crates/infrastructure/tests/}:
\begin{itemize}
  \item \rutarepo{audit_query_postgres.rs}: consulta y autorización.
  \item \rutarepo{audit_query_schema.rs}: migración y esquema.
  \item \rutarepo{audit_query_restore.rs}: respaldo y restauración.
\end{itemize}
Requieren una base
desechable configurada mediante \texttt{CASE\_TEST\_DATABASE\_URL}; sin esa
variable, el retorno anticipado del fixture no demuestra aceptación del almacén.
El último grupo utiliza \texttt{pg\_dump} y \texttt{pg\_restore} sobre un esquema
independiente, además de verificar las columnas originales y sus claves derivadas.\par}

La aceptación integrada aprobada está implementada en
\rutarepo{scripts/api-audit-events-demo.py}, invocado por
\rutarepo{scripts/api-demo.sh}. Contrasta las respuestas con historia y cadena
almacenadas, prueba autorización y continuaciones, y vuelve a comprobarlas tras
restaurar con autenticación renovada. Los dos recorridos de servicios reales en
\rutarepo{web/apps/qadra/tests/live/audit-events.spec.mjs} separan la consulta autorizada de
Owner de la denegación a Litigante. Los resultados y tiempos se registran en
\rutarepo{docs/verification-report.md}; no se atribuyen a una campaña de CI.

La extensión de procedencia conserva sus pruebas de codificación en
\rutarepo{crates/domain/src/audit/context.rs}, las de persistencia en
\rutarepo{crates/infrastructure/tests/audit_context_postgres.rs} y las del
límite HTTP en \rutarepo{crates/web/src/access_audit/}. Distinguen contexto
histórico ausente, identidad admitida, origen observado, cancelación y cierre.
La matriz finita y el protocolo previo se conservan en
\rutarepo{docs/access-audit-matrix.md} y
\rutarepo{docs/access-audit-measurement.md}.

La medición del servicio durable es un caso manual explícito de
\rutarepo{crates/infrastructure/tests/audit_access_postgres.rs}; no se ejecuta
por omisión en CI. El recolector
\rutarepo{scripts/api-access-audit-demo.py} contrasta por UUID los tiempos del
límite HTTP con los eventos confirmados en PostgreSQL. Sus poblaciones y límites
se declaran por separado; los resultados están en
\cref{sec:pruebas-bitacora-contexto}. La restauración conserva tanto los bytes
antiguos como la extensión nueva, sin reconstruir identidad retrospectiva.

Además de los límites propios de la consulta, esta trazabilidad conserva límites
generales del prototipo. La evaluación con personal del despacho requiere
participantes reales y un registro de tareas, terminaciones, errores, asistencia
y observaciones; el kit y sus plantillas están preparados, sin participantes ni
resultados humanos (\cref{sec:pruebas-protocolo-humano}). La TSA local no
acredita una campaña con un proveedor externo; el alcance de ese proveedor sigue
separado conforme a \rutarepo{docs/adr/0009-local-timestamp-authority.md}. En la
cadena de auditoría, la ausencia de anclaje de la cabeza y los campos históricos
de contexto faltantes no se ocultan con el conteo de pruebas.
